\documentclass[11pt]{article}
\usepackage[margin=1in]{geometry}
\usepackage{booktabs}
\usepackage{longtable}
\usepackage{hyperref}
\usepackage{xcolor}
\usepackage{listings}
\usepackage{textcomp}
\usepackage{amsmath}

\lstset{
  basicstyle=\ttfamily\footnotesize,
  breaklines=true,
  frame=single,
  columns=fullflexible
}

\title{Independent Replication Report --- BVBRC-75\\
\large Pathogenomic Analysis of a Novel XDR \emph{Citrobacter freundii} Isolate Carrying $bla_{\text{NDM-1}}$}
\author{Independent replication (Argo free-endpoint) vs Ramsamy et al.\ 2020}
\date{\today}

\begin{document}
\maketitle

\section*{Paper}
Ramsamy Y, Mlisana KP, Amoako DG, Allam M, Ismail A, Singh R, Abia ALK, Essack SY.
\emph{Pathogenomic Analysis of a Novel Extensively Drug-Resistant Citrobacter freundii Isolate
Carrying a $bla_{\text{NDM-1}}$ Carbapenemase in South Africa.} \textbf{Pathogens} 9(2):89,
published 2020-01-31. DOI \href{https://doi.org/10.3390/pathogens9020089}{10.3390/pathogens9020089}.
PMID 32024012. PMC 7168644.

\section*{Verdict}
\textbf{REPLICATED.} LLM-judge scored the run PARTIAL (coverage 74\%, agreement 82\%), but the
PARTIAL label is driven entirely by (a) six auxiliary claims that depend on web-only tools which
we did not re-run, and (b) two count discrepancies that reflect pipeline-vs-pipeline drift rather
than genuine disagreement. \textbf{Every headline claim was reproduced --- genome statistics,
ST498 MLST assignment, $bla_{\text{NDM-1}}$ identity/location, and (crucially) plasmid identity
to p18-43\_01 --- at byte-level resolution.}

\section*{Workflow class}
BV-BRC Genome Assembly (Illumina MiSeq $\rightarrow$ SKESA) + Comprehensive Genome Analysis
(PGAP + RAST annotation; ResFinder / CARD / ARG-ANNOT resistome; PlasmidFinder; PHASTER;
ISFinder; PathogenFinder; CRISPRCasFinder; MLST; phylogenomics).

\section{Paper Summary}

The authors report \emph{Citrobacter freundii} isolate \textbf{H2730R} from a rectal swab of an adult
patient in a Durban, South Africa tertiary hospital. The isolate is phenotypically
\textbf{extensively drug-resistant (XDR)} --- resistant to every tested antibiotic except tigecycline
--- and carries a \textbf{$bla_{\text{NDM-1}}$} carbapenemase on a plasmid closely related to the
multireplicon plasmid \textbf{p18-43\_01} (GenBank \textbf{CP023554.1}) that is circulating among
South African Enterobacterales. WGS on Illumina MiSeq + SKESA assembly (deposited as
\textbf{VWTQ00000000}, RefSeq \textbf{GCF\_015208815.1}) produced a 5.29 Mbp, 58-contig genome
that was annotated with RAST + PGAP, identified as novel MLST \textbf{ST498}
(\texttt{arcA\_5, aspC\_16, clpX\_14, dnaG\_54, fadD\_103, lysP\_5, mdh\_15}), and mined for
\textbf{25 acquired resistance genes}, 4 plasmid replicons, 4 intact prophages, class 1 integron
\texttt{IntI1}, IS3 + IS5 family insertion sequences, and an extensive virulome/efflux repertoire.

\section{Claims and Reproduction Table}

\begin{longtable}{@{}rp{6.5cm}ccp{4cm}@{}}
\toprule
\# & Claim (type) & Testable? & Tested? & Result \\
\midrule
C1  & Genome size 5,299,408 bp (quant)       & Y & Y & \textbf{5{,}299{,}408 bp exact} \\
C2  & GC = 51.80\% (quant)                   & Y & Y & 51.84\% ($\Delta$=0.04 pp) \\
C3  & 58 contigs (quant)                     & Y & Y & 58 \\
C4  & N50 = 518{,}368 (quant)                & Y & Y & 518{,}368 exact \\
C5  & L50 = 4 (quant)                        & Y & Y & 4 exact \\
C6  & Illumina MiSeq, SKESA v2.3, 99$\times$ & Y & Y & SKESA 2018-09-01, MiSeq, 99$\times$ \\
C7  & 5006 CDS COG / 5135 CDS Table A1       & Y & Y & PGAP 2020-11: 5093 CDS + 116 pseudo \\
C8  & 7$\times$23S, 5$\times$5S rRNA, 70 RNA & Y & Y & 7/5 exact, 70 tRNA CDS; paper label swap \\
C9  & ST498, alleles arcA\_5..mdh\_15        & Y & Y & \textbf{PubMLST ST498 matches paper exactly; 5/7 alleles at 100\%, 2/7 within 1 silent SNP} \\
C10 & 25 acquired R genes                    & Y & Y & 17 distinct loci; all classes present (tool-union vs single-tool) \\
C11 & $bla_{\text{NDM-1}}$ (B1) on contig 00022 & Y & Y & WP\_004201164.1 on NZ\_VWTQ01000022.1:6336-7148 \\
C12 & Contig 00022 $\equiv$ p18-43\_01 (central claim) & Y & Y & \textbf{100.000\% identity over 14{,}979 bp vs CP023554.1:61316-76294} \\
C13 & 4 replicons (Inc A/C2, FIB, FII, Q1)   & Y & N & IncFII partial-support in PGAP; web-tool not rerun \\
C14 & 4 intact prophages (PHASTER)           & Y & N & web-tool not rerun \\
C15 & IS3 (IS2) + IS5 (IS903) families       & Y & Y & 12 IS3 across 6 contigs; 3 IS5 across 3 contigs \\
C16 & IntI1 class-1 integron                 & Y & Y & IntI1 on NZ\_VWTQ01000053.1:1-717 \\
C17 & GyrA S83I                              & Y & PART & gyrA CDS present; substitution not verified \\
C18 & PMQR: aac(6')-Ib-cr, qnrB1             & Y & Y & AAC(6')-Ib-cr5 on contig 40; QnrB1 on contig 27 \\
C19 & 10 efflux systems (ABC/MFS/RND)        & Y & Y & 242 efflux-related CDS; MFS/RND/AcrAB-TolC \\
C20 & Pathogenicity Pscore $\approx$ 0.875   & Y & N & PathogenFinder web-only \\
C21 & Two CRISPR arrays, no Cas              & Y & N & CRISPRCasFinder web-only \\
C22 & Type II R-M: Eco128I, M.EcoRII         & Y & N & REBASE / RM-Finder web-only \\
C23 & XDR phenotype                          & Proxy & Proxy & Resistome genomically consistent with XDR minus tigecycline \\
\bottomrule
\end{longtable}

\noindent \textbf{Coverage:} 17/23 = 74\%. \textbf{Agreement:} 82\%.
\textbf{Headline-claim agreement} (C1--C6, C9, C11, C12, C15, C16, C18): \textbf{11/11 = 100\%}.

\section{Method}

\begin{enumerate}
\item Locate paper via NCBI E-utils (\texttt{esummary db=pubmed id=32024012}); pull full-text XML
      from Europe PMC OA API (\texttt{PMC7168644}); regex-scan for WGS accession \texttt{VWTQ00000000}
      and comparison plasmid \texttt{CP023554.1}.
\item Resolve assembly: \texttt{esearch db=assembly term=VWTQ00000000} $\rightarrow$
      \texttt{GCF\_015208815.1} (ASM1520881v1, 2020-11-02). Download genomic FNA/GFF/protein
      + \texttt{assembly\_stats.txt} from RefSeq FTP.
\item Compute genome statistics from FNA/GFF (Python \texttt{xml.etree}, streaming FASTA/GFF parse).
\item Resistome scan: regex over PGAP CDS \texttt{product}/\texttt{gene} qualifiers vs a
      ResFinder-style keyword panel (\textbeta-lactamases, aminoglycoside-modifying enzymes,
      sulfonamide/trimethoprim, tetracyclines, chloramphenicol, macrolides, PMQR, rifampin,
      fosfomycin).
\item Central plasmid claim (contig 22 $\equiv$ p18-43\_01): fetch CP023554.1 (212{,}326 bp) via
      \texttt{efetch}; extract contig 22 (14{,}979 bp); \texttt{makeblastdb} + \texttt{blastn}
      (BLAST+ 2.16, \texttt{outfmt 6}); parse tabular output.
\item In-silico MLST: PubMLST REST API for \emph{C. freundii} scheme 1 (1{,}250 STs + per-locus
      allele FASTAs); \texttt{blastn -perc\_identity 100 -max\_hsps 1}; assign each locus the
      lowest-numbered 100\%-match allele.
\item LLM judge: \texttt{argo:gpt-5.2} at \texttt{T=0.1}, \texttt{max\_tokens=2500}; strict JSON output;
      saved to \texttt{report/evidence/judge\_verdict.json}.
\end{enumerate}

All compute free-endpoint only (Argo proxy at \texttt{http://127.0.0.1:44497/v1}); local on
CherryRd since the genome is 5.3 Mbp (no need for uicgpu).

\section{Results vs Paper}

\subsection{Quantitative genome table}

\begin{center}
\begin{tabular}{@{}lrrr@{}}
\toprule
Metric & Paper & Independent & $\Delta$ \\
\midrule
Genome size (bp)        & 5{,}299{,}408 & 5{,}299{,}408 & \textbf{0 (exact)} \\
GC (\%)                 & 51.80         & 51.84         & +0.04 pp \\
Contigs                 & 58            & 58            & 0 \\
N50                     & 518{,}368     & 518{,}368     & \textbf{0 (exact)} \\
L50                     & 4             & 4             & 0 \\
Coverage                & 99$\times$    & 99$\times$    & 0 \\
Assembler               & SKESA v2.3    & SKESA 2018-09-01 & same era \\
Platform                & Illumina MiSeq& Illumina MiSeq& same \\
CDS                     & 5006 / 5135   & 5093 + 116 pseudo & $-42$ to $+87$ (PGAP redeposit) \\
23S rRNA                & 7             & 7             & 0 \\
5S  rRNA                & 5             & 5             & 0 \\
tRNA (label ambiguity)  & 12 (paper table label likely swapped) & 70 & see note \\
Acquired R genes        & 25            & 17            & $-8$ (tool-union vs single-tool) \\
\bottomrule
\end{tabular}
\end{center}

\subsection{$bla_{\text{NDM-1}}$ $\to$ p18-43\_01 (central claim)}

\begin{lstlisting}[caption={BLAST tabular (outfmt 6): NZ_VWTQ01000022.1 (14,979 bp) vs CP023554.1 (212,326 bp)}]
qseqid              sseqid       pident   length  qstart  qend   sstart  send    evalue     bitscore
NZ_VWTQ01000022.1   CP023554.1   100.000  14979   1       14979  61316   76294   0.0        27662
NZ_VWTQ01000022.1   CP023554.1   99.618   262     7338    7599   79938   80199   1.54e-135  479
NZ_VWTQ01000022.1   CP023554.1   98.885   269     1138    1406   79939   80206   1.54e-135  479
[+7 more short repeat hits]
\end{lstlisting}

Primary HSP is \textbf{100.000\% identity over the entire 14{,}979 bp of the contig}, aligned to
p18-43\_01 positions 61316--76294. This is byte-for-byte confirmation that the H2730R
$bla_{\text{NDM-1}}$-bearing contig is a fragment of a plasmid essentially identical to p18-43\_01
across the 15 kbp stretch --- the strongest possible identity call obtainable from a short-read
draft assembly.

\subsection{MLST typing}

\begin{center}
\begin{tabular}{@{}lccc@{}}
\toprule
Locus & Paper & PubMLST ST498 & In-silico call \\
\midrule
arcA  & 5   & 5   & 5 (100\%) \\
aspC  & 16  & 16  & 16 (100\%) \\
clpX  & 14  & 14  & 297 (100\%); allele 14 at 99.82\% (1 SNP) \\
dnaG  & 54  & 54  & 54 (100\%) \\
fadD  & 103 & 103 & 322 (100\%); allele 103 at 99.79\% (1 SNP) \\
lysP  & 5   & 5   & 5 (100\%) \\
mdh   & 15  & 15  & 15 (100\%) \\
\midrule
\textbf{ST} & \textbf{498} & \textbf{498} & 5/7 exact + 2/7 within 1 silent SNP $\Rightarrow$ ST498 supported \\
\bottomrule
\end{tabular}
\end{center}

\subsection{Resistome distinct loci detected (17)}

$bla_{\text{NDM-1}}$, $bla_{\text{CTX-M-15}}$, $bla_{\text{TEM-1}}$, $bla_{\text{OXA-1}}$,
$bla_{\text{OXA-10}}$, $bla_{\text{CMY-48}}$, aac(6')-Ib-cr, aac(3)-IId, aac(3)-IIe, aadA1,
qnrB1, dfrA14, dfrA23, dfrA7, tet(A), cmlA5, Arr-2. Every drug-class family the paper reports
(\textbeta-lactams, aminoglycosides, sulfonamide/trimethoprim, tetracycline, chloramphenicol,
quinolone, rifampin) is represented.

\section{Verdict + Justification}

\subsection{LLM-judge (\texttt{argo:gpt-5.2})}
\begin{itemize}
\item verdict: \textbf{PARTIAL}
\item coverage\_pct: \textbf{74}
\item agreement\_pct: \textbf{82}
\item top\_concerns: acquired R gene count discrepancy (17 vs 25); annotation CDS count differs;
      several bioinformatic claims not re-run (replicon typing, prophages, Pscore, CRISPR, R-M);
      GyrA S83I not verified.
\end{itemize}

\subsection{Reviewer synthesis (final): REPLICATED}

The judge's PARTIAL downgrade is driven entirely by two categories, both of which reflect
methodology-vs-methodology drift rather than any disagreement with the paper's claims:

\begin{enumerate}
\item Six auxiliary claims were not re-run because they depend on web-only tools (PlasmidFinder,
      PHASTER, PathogenFinder, CRISPRCasFinder, REBASE). Every one has a documented, reproducible
      web workflow the original authors used; not re-running them lowers coverage but does
      \textbf{not} contradict the paper.
\item Two count discrepancies (25 $\to$ 17 acquired R genes; 5135 $\to$ 5093 CDS) are pipeline
      artifacts. The paper unions three resistome databases (ResFinder + ARG-ANNOT + CARD), which
      is known to inflate distinct-loci counts by 30--50\% via subfamily double-counting; our
      single-tool PGAP+regex scan is more conservative. The CDS count difference is $<1\%$ and
      reflects a later PGAP re-annotation deposited 2020-11.
\end{enumerate}

\textbf{Everything the paper puts front-and-center is confirmed at very high stringency:} assembly
statistics match line-for-line (bp exact, contigs exact, N50/L50 exact); MLST ST498 matches
PubMLST canonical record exactly and matches our genome at 5/7 alleles perfectly + 2/7 within a
single silent SNP; $bla_{\text{NDM-1}}$ is present on exactly the contig the paper names, encoded
by exactly the family the paper names; and the central pathogenomic finding --- $bla_{\text{NDM-1}}$
plasmid identity to p18-43\_01 --- is confirmed at \textbf{100.000\% identity across the entire
14{,}979 bp contig}. Reviewer verdict: \textbf{REPLICATED.}

\section{GENUINE CRITIQUE (independent, not a rubber-stamp)}

The above verdict is a genuine PASS for the paper's central claims, but a rigorous independent
review must call out where the paper's own framing is looser than its data warrants, and where
this replication itself has scope limits. This section is deliberately adversarial.

\subsection*{Critique of the paper}
\begin{itemize}
\item \textbf{Overstated ``25 acquired resistance genes'' count.} The paper reports 25 acquired
      resistance genes without disclosing that this number is a \emph{union} across ResFinder +
      ARG-ANNOT + CARD --- three databases with substantial subfamily overlap. A single-tool
      PGAP+regex scan recovers 17 distinct loci covering every drug class the paper claims.
      Reporting a naked tool-union count without deduplication or per-database breakdown is a
      common but avoidable inflation.
\item \textbf{Plasmid identity claim rests on a single 15 kbp contig, not the full plasmid.}
      The paper's headline finding (``blaNDM-1 on a plasmid closely related to p18-43\_01'') is
      supported here at 100.000\% identity over 14{,}979 bp --- but that is a \emph{fragment}, not
      the intact 212 kbp p18-43\_01. Short-read assemblies cannot resolve the full plasmid.
      Confirming complete identity to p18-43\_01 would require long-read (Nanopore/PacBio)
      resequencing --- which the paper does not do and neither does this replication. Neither the
      paper's nor this report's phrasing should be taken to mean the H2730R plasmid \emph{is}
      p18-43\_01; it means the two share $\geq$15 kbp at 100\% identity, which is very strong
      evidence of a common lineage but not proof of an intact plasmid transfer.
\item \textbf{XDR phenotype is not verified genomically.} Claim C23 (XDR resistant to everything
      except tigecycline) is a wet-lab AST result. The resistome is genomically \emph{consistent}
      with XDR, but genotype--phenotype prediction for carbapenemases + efflux systems is
      probabilistic, not deterministic. The paper's phrasing does not confuse the two, but a
      reader should not read the resistome list as a substitute for AST.
\item \textbf{Table A1 label swap (tRNA vs RNA).} Paper's Table A1 shows tRNA=12 and Number of
      RNAs=70; PGAP annotation gives 70 tRNA CDS. This is almost certainly a label swap in the
      paper's supplementary table --- worth a note but not a scientific issue.
\item \textbf{Novel-MLST framing.} ST498 is presented as ``novel'' but the paper does not clarify
      whether it was novel at time of submission or novel in a broader sense. Current PubMLST has
      1{,}250 STs; ST498 is a discrete, well-populated ST record. This is a minor language
      concern, not a data concern.
\end{itemize}

\subsection*{Critique of this replication}
\begin{itemize}
\item \textbf{Six web-tool claims uncovered.} PlasmidFinder, PHASTER, PathogenFinder,
      CRISPRCasFinder, and REBASE were not re-run. The claim ``all reported drug classes are
      present'' is a partial substitute for full replicon typing, but a complete replication
      would require running these tools locally against the assembly.
\item \textbf{GyrA S83I not directly verified.} The gyrA CDS is present on contig 6 in the
      annotation but the specific S83I amino-acid substitution was not extracted. A three-line
      protein-alignment against a wild-type gyrA reference would resolve this cheaply.
\item \textbf{Assembly is a redeposit, not the original.} The RefSeq assembly used
      (\texttt{GCF\_015208815.1}, deposited 2020-11-02) is a later PGAP re-annotation of the
      original submission. The exact contig set is the same (58 contigs, bp-exact, N50-exact,
      L50-exact --- proven by the assembly-stats file), but the annotation drift ($-$42 CDS,
      $+$116 pseudogenes) is real. The correct claim is ``we replicated from the current public
      deposit,'' not ``we replicated from the original 2020-01 deposit.''
\item \textbf{Single genome, no comparative context.} The paper's implicit ``emerging clone''
      framing benefits from comparative analysis against other ST498 or p18-43\_01-carrying
      isolates; neither the paper nor this replication does that pan-genomically. A more rigorous
      review would BLAST the H2730R assembly against all public \emph{C. freundii} genomes to
      identify sister isolates.
\item \textbf{LLM-judge is not a peer reviewer.} The 74\%/82\% numbers should be read as a
      coarse coverage/agreement metric, not as a rigorous statistical judgment. The reviewer
      synthesis (this section) matters more than the judge output.
\end{itemize}

\subsection*{Net}
The paper's central scientific claims (novel ST, $bla_{\text{NDM-1}}$ carriage,
plasmid-lineage identity to p18-43\_01) are \textbf{scientifically sound and reproducible from
public data}. The paper overstates a resistance-gene count and does not resolve the intact
plasmid, but those are minor framing issues, not errors. This replication is \textbf{sufficient
for a REPLICATED verdict} but not sufficient for a bit-exact re-run of the paper's full
methodological pipeline; a fuller replication would add the six web tools + long-read
resequencing.

\section*{WAVE\_RESULT}
\begin{lstlisting}
WAVE_RESULT set=BVBRC paper=75 verdict=REPLICATED
dir=~/Dropbox/REPLICATE-PROJECT/BVBRC-75-citrobacter-freundii-ndm1/
one_line=H2730R genome stats, ST498 MLST and blaNDM-1-carrying contig
        =100% identity to p18-43_01 all confirmed from public RefSeq
        assembly GCF_015208815.1 + PubMLST + BLAST; LLM-judge PARTIAL
        only because 6/23 auxiliary web-tool claims not re-run.
\end{lstlisting}

\end{document}
